\part{CUDA acceleration and performance}

\chapter{From CPU kernels to CUDA kernels}
\label{ch:cuda-from-cpu}
CUDA acceleration begins from a CPU-correct tensor contract. Exact API details should be checked against official CUDA documentation for the installed toolkit. \citep{cudaProgrammingGuide,cudaArchive,cuda13ReleaseNotes}

\begin{figure}[htbp]
\centering
\resizebox{0.94\textwidth}{!}{%
\begin{tikzpicture}[node distance=8mm,>=Latex]
\tikzstyle{level}=[rounded rectangle,draw=EngineBlue!80,fill=EngineBlue!6,minimum width=28mm,minimum height=8mm,align=center,font=\small]
\node[level] (grid) {grid};
\node[level,below left=of grid] (block1) {block};
\node[level,below right=of grid] (block2) {block};
\node[level,below=of block1] (warp1) {warp};
\node[level,below=of block2] (warp2) {warp};
\node[level,below=of warp1] (thread1) {threads};
\node[level,below=of warp2] (thread2) {threads};
\draw[->,thick] (grid)--(block1);
\draw[->,thick] (grid)--(block2);
\draw[->,thick] (block1)--(warp1);
\draw[->,thick] (block2)--(warp2);
\draw[->,thick] (warp1)--(thread1);
\draw[->,thick] (warp2)--(thread2);
\node[rounded rectangle,draw=EngineGreen!80,fill=EngineGreen!7,align=left,text width=44mm,right=15mm of block2,font=\small] {registers\\shared memory\\global memory\\constant memory};
\end{tikzpicture}
}
\caption{CUDA execution and memory hierarchy.}
\end{figure}

\begin{table}[htbp]\centering\small
\begin{tabularx}{\textwidth}{L{0.20\textwidth}Y}
\toprule
Concept & Engine meaning \\
\midrule
Thread & Scalar work lane. \\
Warp & Scheduling group; coalescing and divergence matter. \\
Block & Cooperative group with shared memory. \\
Stream & Ordered queue of GPU work. \\
Event & Timing and dependency marker. \\
CUDA Graph & Captured repeated launch graph. \\
\bottomrule
\end{tabularx}
\caption{CUDA execution checklist.}
\end{table}

\begin{intuitionbox}{Intuition}
A GPU kernel is a parallel implementation of an already proven operation.
\end{intuitionbox}

\begin{definitionbox}{Mathematical or contract statement}
CUDA executes grids of blocks of threads, scheduled in warps.
\end{definitionbox}

\begin{implementationbox}{Implementation interpretation}
Compare every CUDA output against CPU/PyTorch golden vectors before profiling.
\end{implementationbox}

\begin{verificationbox}{How to verify this works}
\begin{itemize}
\item Check shapes and dtypes at the boundary.
\item Compare deterministic outputs against PyTorch golden vectors.
\item Record tolerance, seed, model artifact, and backend.
\item Do not benchmark until validation passes.
\end{itemize}
\end{verificationbox}

\begin{warningbox}{Common failure modes}
\begin{itemize}
\item Letting framework defaults choose dtype or layout.
\item Testing only a batch size of one.
\item Depending on upstream checkpoint names in runtime code.
\item Treating benchmark output as validation.
\end{itemize}
\end{warningbox}

\begin{chapterexercises}
\item State the central invariant in one sentence.\\
\textit{Hint.} Use the conceptual guide vocabulary.
\item Construct a toy example with small shapes.\\
\textit{Hint.} Use $B=2$, $T=3$, $H=8$ where possible.
\item Name a validation test that would catch a plausible bug.\\
\textit{Hint.} Prefer generated golden vectors to handwritten long arrays.
\end{chapterexercises}

\chapter{Memory layout and tensor storage}
\label{ch:gpu-memory-layout}
GPU performance is often dominated by memory traffic, so layout decisions matter.

\begin{table}[htbp]\centering\small
\begin{tabularx}{\textwidth}{L{0.22\textwidth}Y}
\toprule
Memory & Engine concern \\
\midrule
Registers & Per-thread temporaries and spill risk. \\
Shared memory & Block-level staging and reductions. \\
Global memory & DRAM bandwidth and coalescing. \\
Constant memory & Small broadcast-friendly values. \\
Host memory & Transfers, pinning, and orchestration. \\
\bottomrule
\end{tabularx}
\caption{Memory hierarchy as a design checklist.}
\end{table}

\begin{intuitionbox}{Intuition}
Layout is where algebra becomes bandwidth.
\end{intuitionbox}

\begin{definitionbox}{Mathematical or contract statement}
Artifact row-major layout is a serialisation contract; GPU internal layout may differ if logical indexing is preserved.
\end{definitionbox}

\begin{implementationbox}{Implementation interpretation}
Repack weights only after validating equality with artifact tensors.
\end{implementationbox}

\begin{verificationbox}{How to verify this works}
\begin{itemize}
\item Check shapes and dtypes at the boundary.
\item Compare deterministic outputs against PyTorch golden vectors.
\item Record tolerance, seed, model artifact, and backend.
\item Do not benchmark until validation passes.
\end{itemize}
\end{verificationbox}

\begin{warningbox}{Common failure modes}
\begin{itemize}
\item Letting framework defaults choose dtype or layout.
\item Testing only a batch size of one.
\item Depending on upstream checkpoint names in runtime code.
\item Treating benchmark output as validation.
\end{itemize}
\end{warningbox}

\begin{chapterexercises}
\item State the central invariant in one sentence.\\
\textit{Hint.} Use the conceptual guide vocabulary.
\item Construct a toy example with small shapes.\\
\textit{Hint.} Use $B=2$, $T=3$, $H=8$ where possible.
\item Name a validation test that would catch a plausible bug.\\
\textit{Hint.} Prefer generated golden vectors to handwritten long arrays.
\end{chapterexercises}

\chapter{CUDA matmul concepts}
\label{ch:cuda-matmul}
Matmul dominates projection and MLP work. CUTLASS-style ideas use tiling, shared memory, and Tensor Cores. \citep{cutlassDocs}

\[
\text{arithmetic intensity}\approx \frac{2MNK}{\text{bytes read and written}}.
\]

\begin{intuitionbox}{Intuition}
Matmul speed comes from reusing data after reading it.
\end{intuitionbox}

\begin{definitionbox}{Mathematical or contract statement}
For \(A\in\R^{M\times K}\) and \(B\in\R^{K\times N}\), computing \(C=AB\) performs \(MNK\) multiply-adds, commonly counted as roughly \(2MNK\) floating-point operations.
\end{definitionbox}

\begin{implementationbox}{Implementation interpretation}
Prefer trusted libraries for production while using toy kernels for education.
\end{implementationbox}

\begin{verificationbox}{How to verify this works}
\begin{itemize}
\item Check shapes and dtypes at the boundary.
\item Compare deterministic outputs against PyTorch golden vectors.
\item Record tolerance, seed, model artifact, and backend.
\item Do not benchmark until validation passes.
\end{itemize}
\end{verificationbox}

\begin{warningbox}{Common failure modes}
\begin{itemize}
\item Letting framework defaults choose dtype or layout.
\item Testing only a batch size of one.
\item Depending on upstream checkpoint names in runtime code.
\item Treating benchmark output as validation.
\end{itemize}
\end{warningbox}

\begin{chapterexercises}
\item State the central invariant in one sentence.\\
\textit{Hint.} Use the conceptual guide vocabulary.
\item Construct a toy example with small shapes.\\
\textit{Hint.} Use $B=2$, $T=3$, $H=8$ where possible.
\item Name a validation test that would catch a plausible bug.\\
\textit{Hint.} Prefer generated golden vectors to handwritten long arrays.
\end{chapterexercises}

\chapter{LayerNorm and RMSNorm kernels}
\label{ch:norm-kernels}
Normalisation kernels are row-wise reductions over hidden size.



\begin{intuitionbox}{Intuition}
A norm kernel is a reduction plus an elementwise transform.
\end{intuitionbox}

\begin{definitionbox}{Mathematical or contract statement}
LayerNorm reduces mean and variance; RMSNorm reduces mean square.
\end{definitionbox}

\begin{implementationbox}{Implementation interpretation}
Map rows to warps or blocks and validate epsilon and axis semantics.
\end{implementationbox}

\begin{verificationbox}{How to verify this works}
\begin{itemize}
\item Check shapes and dtypes at the boundary.
\item Compare deterministic outputs against PyTorch golden vectors.
\item Record tolerance, seed, model artifact, and backend.
\item Do not benchmark until validation passes.
\end{itemize}
\end{verificationbox}

\begin{warningbox}{Common failure modes}
\begin{itemize}
\item Letting framework defaults choose dtype or layout.
\item Testing only a batch size of one.
\item Depending on upstream checkpoint names in runtime code.
\item Treating benchmark output as validation.
\end{itemize}
\end{warningbox}

\begin{chapterexercises}
\item State the central invariant in one sentence.\\
\textit{Hint.} Use the conceptual guide vocabulary.
\item Construct a toy example with small shapes.\\
\textit{Hint.} Use $B=2$, $T=3$, $H=8$ where possible.
\item Name a validation test that would catch a plausible bug.\\
\textit{Hint.} Prefer generated golden vectors to handwritten long arrays.
\end{chapterexercises}

\chapter{Softmax kernels}
\label{ch:softmax-kernels}
Softmax appears in attention and sampling.

A stable softmax kernel performs a maximum reduction, an exponential sum reduction after max subtraction, and normalisation. Fused attention avoids materialising large probability matrices. \citep{dao2022flashattention,dao2024flashattention2}

\begin{intuitionbox}{Intuition}
Stable softmax is a reduction pipeline.
\end{intuitionbox}

\begin{definitionbox}{Mathematical or contract statement}
Subtract the maximum, sum exponentials, and normalise, applying masks before normalisation.
\end{definitionbox}

\begin{implementationbox}{Implementation interpretation}
Validate shift invariance and masked support on GPU.
\end{implementationbox}

\begin{verificationbox}{How to verify this works}
\begin{itemize}
\item Check shapes and dtypes at the boundary.
\item Compare deterministic outputs against PyTorch golden vectors.
\item Record tolerance, seed, model artifact, and backend.
\item Do not benchmark until validation passes.
\end{itemize}
\end{verificationbox}

\begin{warningbox}{Common failure modes}
\begin{itemize}
\item Letting framework defaults choose dtype or layout.
\item Testing only a batch size of one.
\item Depending on upstream checkpoint names in runtime code.
\item Treating benchmark output as validation.
\end{itemize}
\end{warningbox}

\begin{chapterexercises}
\item State the central invariant in one sentence.\\
\textit{Hint.} Use the conceptual guide vocabulary.
\item Construct a toy example with small shapes.\\
\textit{Hint.} Use $B=2$, $T=3$, $H=8$ where possible.
\item Name a validation test that would catch a plausible bug.\\
\textit{Hint.} Prefer generated golden vectors to handwritten long arrays.
\end{chapterexercises}

\chapter{Attention kernels}
\label{ch:attention-kernels}
Attention kernels combine score computation, masking, softmax, and value mixing. \citep{dao2022flashattention,dao2024flashattention2}



\begin{intuitionbox}{Intuition}
The costly object is often the intermediate matrix.
\end{intuitionbox}

\begin{definitionbox}{Mathematical or contract statement}
Exact attention can stream tiles while maintaining stable softmax statistics.
\end{definitionbox}

\begin{implementationbox}{Implementation interpretation}
Keep unfused and fused paths comparable.
\end{implementationbox}

\begin{verificationbox}{How to verify this works}
\begin{itemize}
\item Check shapes and dtypes at the boundary.
\item Compare deterministic outputs against PyTorch golden vectors.
\item Record tolerance, seed, model artifact, and backend.
\item Do not benchmark until validation passes.
\end{itemize}
\end{verificationbox}

\begin{warningbox}{Common failure modes}
\begin{itemize}
\item Letting framework defaults choose dtype or layout.
\item Testing only a batch size of one.
\item Depending on upstream checkpoint names in runtime code.
\item Treating benchmark output as validation.
\end{itemize}
\end{warningbox}

\begin{chapterexercises}
\item State the central invariant in one sentence.\\
\textit{Hint.} Use the conceptual guide vocabulary.
\item Construct a toy example with small shapes.\\
\textit{Hint.} Use $B=2$, $T=3$, $H=8$ where possible.
\item Name a validation test that would catch a plausible bug.\\
\textit{Hint.} Prefer generated golden vectors to handwritten long arrays.
\end{chapterexercises}

\chapter{Prefill kernels}
\label{ch:prefill-kernels}
Prefill processes prompt sequences with enough parallel work to keep the GPU busy.



\begin{intuitionbox}{Intuition}
Prefill is throughput-oriented.
\end{intuitionbox}

\begin{definitionbox}{Mathematical or contract statement}
Prefill writes cache entries for all prompt positions.
\end{definitionbox}

\begin{implementationbox}{Implementation interpretation}
Benchmark prefill separately from decode.
\end{implementationbox}

\begin{verificationbox}{How to verify this works}
\begin{itemize}
\item Check shapes and dtypes at the boundary.
\item Compare deterministic outputs against PyTorch golden vectors.
\item Record tolerance, seed, model artifact, and backend.
\item Do not benchmark until validation passes.
\end{itemize}
\end{verificationbox}

\begin{warningbox}{Common failure modes}
\begin{itemize}
\item Letting framework defaults choose dtype or layout.
\item Testing only a batch size of one.
\item Depending on upstream checkpoint names in runtime code.
\item Treating benchmark output as validation.
\end{itemize}
\end{warningbox}

\begin{chapterexercises}
\item State the central invariant in one sentence.\\
\textit{Hint.} Use the conceptual guide vocabulary.
\item Construct a toy example with small shapes.\\
\textit{Hint.} Use $B=2$, $T=3$, $H=8$ where possible.
\item Name a validation test that would catch a plausible bug.\\
\textit{Hint.} Prefer generated golden vectors to handwritten long arrays.
\end{chapterexercises}

\chapter{Decode kernels}
\label{ch:decode-kernels}
Decode often processes one token and reads the accumulated KV cache.



\begin{intuitionbox}{Intuition}
Decode is latency and bandwidth sensitive.
\end{intuitionbox}

\begin{definitionbox}{Mathematical or contract statement}
For prompt length \(m\) and zero-based decode step \(t\), each layer reads visible length \(S=m+t\) before appending and writes one new key/value position, giving cache length \(S+1\) after the step.
\end{definitionbox}

\begin{implementationbox}{Implementation interpretation}
Use decode-specific kernels and validate cache equivalence.
\end{implementationbox}

\begin{verificationbox}{How to verify this works}
\begin{itemize}
\item Check shapes and dtypes at the boundary.
\item Compare deterministic outputs against PyTorch golden vectors.
\item Record tolerance, seed, model artifact, and backend.
\item Do not benchmark until validation passes.
\end{itemize}
\end{verificationbox}

\begin{warningbox}{Common failure modes}
\begin{itemize}
\item Letting framework defaults choose dtype or layout.
\item Testing only a batch size of one.
\item Depending on upstream checkpoint names in runtime code.
\item Treating benchmark output as validation.
\end{itemize}
\end{warningbox}

\begin{chapterexercises}
\item State the central invariant in one sentence.\\
\textit{Hint.} Use the conceptual guide vocabulary.
\item Construct a toy example with small shapes.\\
\textit{Hint.} Use $B=2$, $T=3$, $H=8$ where possible.
\item Name a validation test that would catch a plausible bug.\\
\textit{Hint.} Prefer generated golden vectors to handwritten long arrays.
\end{chapterexercises}

\chapter{KV-cache layout on GPU}
\label{ch:gpu-cache-layout}
The GPU KV cache must preserve logical layout while optimising physical memory reads.



\begin{intuitionbox}{Intuition}
The cache is usually the largest live decode state.
\end{intuitionbox}

\begin{definitionbox}{Mathematical or contract statement}
MHA, MQA, and GQA differ primarily in \(H_{kv}\), which controls cache bytes. For element size \(e\) bytes and \(L\) layers, an append-only KV cache stores approximately \(2LBH_{kv}Sd_h e\) bytes.
\end{definitionbox}

\begin{implementationbox}{Implementation interpretation}
Keep logical cache API separate from physical allocation.
\end{implementationbox}

\begin{verificationbox}{How to verify this works}
\begin{itemize}
\item Check shapes and dtypes at the boundary.
\item Compare deterministic outputs against PyTorch golden vectors.
\item Record tolerance, seed, model artifact, and backend.
\item Do not benchmark until validation passes.
\end{itemize}
\end{verificationbox}

\begin{warningbox}{Common failure modes}
\begin{itemize}
\item Letting framework defaults choose dtype or layout.
\item Testing only a batch size of one.
\item Depending on upstream checkpoint names in runtime code.
\item Treating benchmark output as validation.
\end{itemize}
\end{warningbox}

\begin{chapterexercises}
\item State the central invariant in one sentence.\\
\textit{Hint.} Use the conceptual guide vocabulary.
\item Construct a toy example with small shapes.\\
\textit{Hint.} Use $B=2$, $T=3$, $H=8$ where possible.
\item Name a validation test that would catch a plausible bug.\\
\textit{Hint.} Prefer generated golden vectors to handwritten long arrays.
\end{chapterexercises}

\chapter{Kernel fusion}
\label{ch:fusion}
Kernel fusion reduces memory traffic and launch overhead by combining operations.



\begin{intuitionbox}{Intuition}
Fusion should remove memory movement, not validation boundaries.
\end{intuitionbox}

\begin{definitionbox}{Mathematical or contract statement}
Examples include residual plus norm, bias plus activation, and fused attention.
\end{definitionbox}

\begin{implementationbox}{Implementation interpretation}
Keep unfused references for debugging.
\end{implementationbox}

\begin{verificationbox}{How to verify this works}
\begin{itemize}
\item Check shapes and dtypes at the boundary.
\item Compare deterministic outputs against PyTorch golden vectors.
\item Record tolerance, seed, model artifact, and backend.
\item Do not benchmark until validation passes.
\end{itemize}
\end{verificationbox}

\begin{warningbox}{Common failure modes}
\begin{itemize}
\item Letting framework defaults choose dtype or layout.
\item Testing only a batch size of one.
\item Depending on upstream checkpoint names in runtime code.
\item Treating benchmark output as validation.
\end{itemize}
\end{warningbox}

\begin{chapterexercises}
\item State the central invariant in one sentence.\\
\textit{Hint.} Use the conceptual guide vocabulary.
\item Construct a toy example with small shapes.\\
\textit{Hint.} Use $B=2$, $T=3$, $H=8$ where possible.
\item Name a validation test that would catch a plausible bug.\\
\textit{Hint.} Prefer generated golden vectors to handwritten long arrays.
\end{chapterexercises}

\chapter{Tensor Cores and mixed precision}
\label{ch:tensor-cores}
Tensor Cores accelerate supported matrix operations and dtypes. \citep{cudaProgrammingGuide,cutlassDocs}



\begin{intuitionbox}{Intuition}
Specialised hardware is useful only when the numerical contract permits it.
\end{intuitionbox}

\begin{definitionbox}{Mathematical or contract statement}
Float16 and bfloat16 reduce bandwidth but change tolerance expectations.
\end{definitionbox}

\begin{implementationbox}{Implementation interpretation}
Keep phase-0 float32 artifact support separate from later mixed-precision execution.
\end{implementationbox}

\begin{verificationbox}{How to verify this works}
\begin{itemize}
\item Check shapes and dtypes at the boundary.
\item Compare deterministic outputs against PyTorch golden vectors.
\item Record tolerance, seed, model artifact, and backend.
\item Do not benchmark until validation passes.
\end{itemize}
\end{verificationbox}

\begin{warningbox}{Common failure modes}
\begin{itemize}
\item Letting framework defaults choose dtype or layout.
\item Testing only a batch size of one.
\item Depending on upstream checkpoint names in runtime code.
\item Treating benchmark output as validation.
\end{itemize}
\end{warningbox}

\begin{chapterexercises}
\item State the central invariant in one sentence.\\
\textit{Hint.} Use the conceptual guide vocabulary.
\item Construct a toy example with small shapes.\\
\textit{Hint.} Use $B=2$, $T=3$, $H=8$ where possible.
\item Name a validation test that would catch a plausible bug.\\
\textit{Hint.} Prefer generated golden vectors to handwritten long arrays.
\end{chapterexercises}

\chapter{CUDA Graphs and launch overhead}
\label{ch:cuda-graphs}
CUDA Graphs are relevant to repeated decode schedules. \citep{cudaGraphsGuide}

CUDA Graphs can reduce repeated launch overhead, but only after output equality and shape/memory stability are validated. \citep{cudaGraphsGuide}

\begin{intuitionbox}{Intuition}
A repeated schedule can be captured and replayed.
\end{intuitionbox}

\begin{definitionbox}{Mathematical or contract statement}
A graph captures operations and dependencies, reducing launch overhead if shapes and memory addresses are stable.
\end{definitionbox}

\begin{implementationbox}{Implementation interpretation}
Validate non-graph and graph outputs before trusting speed.
\end{implementationbox}

\begin{verificationbox}{How to verify this works}
\begin{itemize}
\item Check shapes and dtypes at the boundary.
\item Compare deterministic outputs against PyTorch golden vectors.
\item Record tolerance, seed, model artifact, and backend.
\item Do not benchmark until validation passes.
\end{itemize}
\end{verificationbox}

\begin{warningbox}{Common failure modes}
\begin{itemize}
\item Letting framework defaults choose dtype or layout.
\item Testing only a batch size of one.
\item Depending on upstream checkpoint names in runtime code.
\item Treating benchmark output as validation.
\end{itemize}
\end{warningbox}

\begin{chapterexercises}
\item State the central invariant in one sentence.\\
\textit{Hint.} Use the conceptual guide vocabulary.
\item Construct a toy example with small shapes.\\
\textit{Hint.} Use $B=2$, $T=3$, $H=8$ where possible.
\item Name a validation test that would catch a plausible bug.\\
\textit{Hint.} Prefer generated golden vectors to handwritten long arrays.
\end{chapterexercises}

\chapter{Profiling workflow}
\label{ch:profiling}
Profiling turns performance guesses into measurements. Nsight Compute and CUDA events are common tools. \citep{nsightCompute,cudaProgrammingGuide}



\begin{intuitionbox}{Intuition}
A profiler explains why, not just how fast.
\end{intuitionbox}

\begin{definitionbox}{Mathematical or contract statement}
Useful diagnostics include memory throughput, warp stalls, Tensor Core utilisation, occupancy, and launch overhead.
\end{definitionbox}

\begin{implementationbox}{Implementation interpretation}
Profile only after correctness passes.
\end{implementationbox}

\begin{verificationbox}{How to verify this works}
\begin{itemize}
\item Check shapes and dtypes at the boundary.
\item Compare deterministic outputs against PyTorch golden vectors.
\item Record tolerance, seed, model artifact, and backend.
\item Do not benchmark until validation passes.
\end{itemize}
\end{verificationbox}

\begin{warningbox}{Common failure modes}
\begin{itemize}
\item Letting framework defaults choose dtype or layout.
\item Testing only a batch size of one.
\item Depending on upstream checkpoint names in runtime code.
\item Treating benchmark output as validation.
\end{itemize}
\end{warningbox}

\begin{chapterexercises}
\item State the central invariant in one sentence.\\
\textit{Hint.} Use the conceptual guide vocabulary.
\item Construct a toy example with small shapes.\\
\textit{Hint.} Use $B=2$, $T=3$, $H=8$ where possible.
\item Name a validation test that would catch a plausible bug.\\
\textit{Hint.} Prefer generated golden vectors to handwritten long arrays.
\end{chapterexercises}

\chapter{Benchmark interpretation}
\label{ch:benchmark-interpretation}
Benchmark interpretation separates prefill throughput, decode latency, first-token latency, and end-to-end generation.



\begin{intuitionbox}{Intuition}
A benchmark number needs a workload.
\end{intuitionbox}

\begin{definitionbox}{Mathematical or contract statement}
Throughput, latency, bandwidth, and arithmetic intensity measure different behaviours.
\end{definitionbox}

\begin{implementationbox}{Implementation interpretation}
Report prompt length, decode length, batch size, backend, and validation status.
\end{implementationbox}

\begin{verificationbox}{How to verify this works}
\begin{itemize}
\item Check shapes and dtypes at the boundary.
\item Compare deterministic outputs against PyTorch golden vectors.
\item Record tolerance, seed, model artifact, and backend.
\item Do not benchmark until validation passes.
\end{itemize}
\end{verificationbox}

\begin{warningbox}{Common failure modes}
\begin{itemize}
\item Letting framework defaults choose dtype or layout.
\item Testing only a batch size of one.
\item Depending on upstream checkpoint names in runtime code.
\item Treating benchmark output as validation.
\end{itemize}
\end{warningbox}

\begin{chapterexercises}
\item State the central invariant in one sentence.\\
\textit{Hint.} Use the conceptual guide vocabulary.
\item Construct a toy example with small shapes.\\
\textit{Hint.} Use $B=2$, $T=3$, $H=8$ where possible.
\item Name a validation test that would catch a plausible bug.\\
\textit{Hint.} Prefer generated golden vectors to handwritten long arrays.
\end{chapterexercises}

\chapter{Optional architecture notes for RTX 40-series and RTX 5060-style GPUs}
\label{ch:rtx-notes}
The target is generic NVIDIA CUDA GPUs; consumer GPUs are useful development targets but not design foundations. \citep{nvidiaRtx40,nvidiaRtx5060,cudaProgrammingGuide}



\begin{intuitionbox}{Intuition}
Portable performance starts from capability queries.
\end{intuitionbox}

\begin{definitionbox}{Mathematical or contract statement}
Query compute capability, memory size, dtype support, shared memory limits, and graph features at runtime.
\end{definitionbox}

\begin{implementationbox}{Implementation interpretation}
Use architecture-specific paths only behind fallbacks.
\end{implementationbox}

\begin{verificationbox}{How to verify this works}
\begin{itemize}
\item Check shapes and dtypes at the boundary.
\item Compare deterministic outputs against PyTorch golden vectors.
\item Record tolerance, seed, model artifact, and backend.
\item Do not benchmark until validation passes.
\end{itemize}
\end{verificationbox}

\begin{warningbox}{Common failure modes}
\begin{itemize}
\item Letting framework defaults choose dtype or layout.
\item Testing only a batch size of one.
\item Depending on upstream checkpoint names in runtime code.
\item Treating benchmark output as validation.
\end{itemize}
\end{warningbox}

\begin{chapterexercises}
\item State the central invariant in one sentence.\\
\textit{Hint.} Use the conceptual guide vocabulary.
\item Construct a toy example with small shapes.\\
\textit{Hint.} Use $B=2$, $T=3$, $H=8$ where possible.
\item Name a validation test that would catch a plausible bug.\\
\textit{Hint.} Prefer generated golden vectors to handwritten long arrays.
\end{chapterexercises}
